\documentclass{amsart}
\usepackage{amsmath,amssymb,amsthm,hyperref}
\usepackage{mathtools}

\newtheorem{theorem}{Theorem}
\newtheorem{lemma}{Lemma}
\newtheorem{proposition}{Proposition}
\newtheorem{corollary}{Corollary}
\theoremstyle{definition}
\newtheorem{definition}{Definition}
\newtheorem{remark}{Remark}

\DeclareMathOperator{\rank}{rank}
\DeclareMathOperator{\range}{range}
\DeclareMathOperator{\tr}{tr}
\DeclareMathOperator{\spann}{span}
\newcommand{\R}{\mathbb{R}}
\newcommand{\Frob}[1]{\left\lVert #1\right\rVert_F}
\newcommand{\Op}[1]{\left\lVert #1\right\rVert_2}
\newcommand{\Atil}{\widetilde{A}}
\newcommand{\E}{\mathbb{E}}
\newcommand{\Prob}{\mathbb{P}}

\begin{document}

\title{Sharp thresholds, failure regimes, and finite-precision limits for
randomized low-rank approximation}
\maketitle

\begin{abstract}
We give a rigorous, self-contained answer to the three questions. For Gaussian
sketching we prove a \emph{matching} two-sided bound: with sketch width
$\ell=k+p$,
\[
\E\,\Frob{A-P_YA}^2\le\Big(1+\tfrac{k}{p-1}\Big)\Frob{A-A_k}^2
\quad(\text{Thm.~\ref{thm:upper}}),
\]
and there is an explicit matrix family attaining this factor in the limit
(Thm.~\ref{thm:lower}); hence the operating regime $p=\Theta(k/\varepsilon)$ is
sharp for the $(1+\varepsilon)$ relative-error guarantee. For oblivious
entrywise sampling we prove a sharp coherence barrier: uniform sampling fails
catastrophically on coherent/sparse matrices unless essentially all entries are
kept, and the general barrier scales as $\mu(A)\,\mathrm{sr}(A)\max(m,n)$
(Thm.~\ref{thm:fail}), removable only by norm-aware (leverage) sampling. Finally
we delineate the finite-precision boundary: the randomized analysis dominates
the deterministic roundoff exactly when the unit roundoff satisfies
$u\,\kappa_{\mathrm{eff}}\Frob{A}\lesssim\varepsilon\,\Frob{A-A_k}$, with
$\kappa_{\mathrm{eff}}=\sigma_1/(\sigma_k-\sigma_{k+1})$ (Thm.~\ref{thm:precision}).
\end{abstract}

\section{Setup and conventions}\label{sec:setup}

Throughout $A\in\R^{m\times n}$ has singular value decomposition (SVD)
$A=\sum_{j=1}^{r}\sigma_j u_j v_j^{\top}$ with
$\sigma_1\ge\cdots\ge\sigma_r>0$, $r=\rank(A)$, and orthonormal
$\{u_j\}\subset\R^m$, $\{v_j\}\subset\R^n$. We write $\Op{\cdot}$ for the
spectral norm and $\Frob{\cdot}$ for the Frobenius norm, so $\Op{A}=\sigma_1$,
$\Frob{A}^2=\sum_j\sigma_j^2$. For a target rank $1\le k\le r$ set
\begin{equation}\label{eq:partition}
A_k=\sum_{j\le k}\sigma_j u_j v_j^\top,\qquad
\tau_k^2:=\Frob{A-A_k}^2=\sum_{j>k}\sigma_j^2 .
\end{equation}
By the Eckart--Young--Mirsky theorem $A_k$ minimizes $\Frob{A-B}$ and
$\Op{A-B}$ over $\rank(B)\le k$, with $\Frob{A-A_k}^2=\tau_k^2$ and
$\Op{A-A_k}=\sigma_{k+1}$, and more generally
\begin{equation}\label{eq:EY}
\Frob{A-B}^2\ \ge\ \sum_{j>\ell}\sigma_j^2\qquad\text{for every }B\text{ with }\rank(B)\le\ell.
\end{equation}

\begin{definition}[Stable rank]\label{def:srank}
$\mathrm{sr}(A):=\Frob{A}^2/\Op{A}^2\in[1,r]$.
\end{definition}
\begin{definition}[Coherence]\label{def:coh}
Let $U\in\R^{m\times r}$, $V\in\R^{n\times r}$ hold the left/right singular
vectors. The coherence is
$\mu(A):=\max\big(\tfrac{m}{r}\max_i\|U_{i,:}\|_2^2,\
\tfrac{n}{r}\max_j\|V_{j,:}\|_2^2\big)\in[1,\max(m,n)/r]$.
\end{definition}

We write $g\lesssim h$ for $g\le Ch$ with an absolute constant $C$. The two
paradigms analyzed are:

\smallskip
\noindent\textbf{(I) Linear sketching + truncation.} Draw
$\Omega\in\R^{n\times\ell}$, form $Y=A\Omega$, let $Q\in\R^{m\times\ell}$ have
orthonormal columns spanning $\range(Y)$, write $P_Y=QQ^\top$, and set
\begin{equation}\label{eq:rand}
A_k^{\mathrm{rand}}:=(P_YA)_k=(QQ^\top A)_k .
\end{equation}
The Gaussian sketch takes $\Omega$ with i.i.d.\ $N(0,1)$ entries.

\smallskip
\noindent\textbf{(II) Entrywise sparsification.} Given $p=(p_{ij})$ and budget
$s$, keep entries independently and rescale,
\begin{equation}\label{eq:entrywise}
\Atil_{ij}=\tfrac{A_{ij}}{q_{ij}}\,\xi_{ij},\quad
\xi_{ij}\sim\mathrm{Bernoulli}(q_{ij}),\ q_{ij}=\min(1,s\,p_{ij}),
\end{equation}
so $\E\Atil=A$, expected support $\le s$, and $A_k^{\mathrm{rand}}=\Atil_k$.

\section{Part (a): sharp sketch-size threshold (Gaussian sketch)}

\subsection{Notation in the singular basis}
Let $V_1\in\R^{n\times k}$, $V_2\in\R^{n\times(r-k)}$ hold $v_1,\dots,v_k$ and
$v_{k+1},\dots,v_r$; let $U_1\in\R^{m\times k}$, $U_2\in\R^{m\times(r-k)}$ hold
$u_1,\dots,u_k$ and $u_{k+1},\dots,u_r$; and
$\Sigma_1=\mathrm{diag}(\sigma_1,\dots,\sigma_k)$,
$\Sigma_2=\mathrm{diag}(\sigma_{k+1},\dots,\sigma_r)$. Write
\begin{equation}\label{eq:OmBlocks}
\Omega_1:=V_1^\top\Omega\in\R^{k\times\ell},\qquad
\Omega_2:=V_2^\top\Omega\in\R^{(r-k)\times\ell}.
\end{equation}
When $\Omega_1$ has full row rank $k$ we write
$\Omega_1^{\dagger}=\Omega_1^\top(\Omega_1\Omega_1^\top)^{-1}\in\R^{\ell\times k}$.

\subsection{Deterministic structural bound}

\begin{lemma}[Structural upper bound]\label{lem:struct}
If $\Omega_1$ has full row rank $k$, then
\begin{equation}\label{eq:structbound}
\Frob{A-P_Y A}^2 \;\le\; \Frob{\Sigma_2}^2
+\Frob{\Sigma_2\,\Omega_2\,\Omega_1^{\dagger}}^2
=\tau_k^2+\Frob{\Sigma_2\,\Omega_2\,\Omega_1^{\dagger}}^2 .
\end{equation}
\end{lemma}

\begin{proof}
The Frobenius norm, the column space of $A\Omega$, and the projector $P_Y$ are
unchanged if we replace $A,\Omega$ by $U^\top A V,\ V^\top\Omega$ (left/right
multiplication by orthogonal $U,V$ permutes coordinates and the projector
transforms accordingly while $\Frob{\cdot}$ is invariant). Hence we may assume
\[
A=\begin{pmatrix}\Sigma_1&0\\0&\Sigma_2\end{pmatrix}\ (\text{padded by zeros}),
\quad U=I,\ V=I,
\]
so $\Omega_1,\Omega_2$ are the top-$k$ and bottom-$(r-k)$ row blocks of
$\Omega$. The product $Y\Omega_1^{\dagger}=A\Omega\,\Omega_1^{\dagger}$ has
columns in $\range(Y)$, and
\[
Y\,\Omega_1^{\dagger}=\Sigma\,\Omega\,\Omega_1^{\dagger}
=\begin{pmatrix}\Sigma_1\,\Omega_1\Omega_1^{\dagger}\\
\Sigma_2\,\Omega_2\Omega_1^{\dagger}\end{pmatrix}
=\begin{pmatrix}\Sigma_1\\ \Sigma_2\,\Omega_2\Omega_1^{\dagger}\end{pmatrix},
\]
using $\Omega_1\Omega_1^{\dagger}=I_k$. Define $W\in\R^{m\times n}$ whose first
$k$ columns equal $Y\Omega_1^{\dagger}$ and whose remaining columns are zero;
then every column of $W$ lies in $\range(Y)$. For each column $a$ of $A$,
$\|a-P_Y a\|_2=\min_{w\in\range(Y)}\|a-w\|_2\le\|a-w_a\|_2$ where $w_a$ is the
corresponding column of $W$; summing squares over columns,
\begin{equation}\label{eq:projmin}
\Frob{A-P_Y A}^2\le\Frob{A-W}^2 .
\end{equation}
Now $A=\mathrm{diag}(\Sigma_1,\Sigma_2)$ has first $k$ columns
$\big(\Sigma_1;0\big)$ and remaining columns $\big(0;\Sigma_2\big)$. On the
first $k$ columns,
$A-W=\big(0;-\Sigma_2\Omega_2\Omega_1^{\dagger}\big)$, contributing
$\Frob{\Sigma_2\Omega_2\Omega_1^{\dagger}}^2$; on the remaining columns
$W=0$, so $A-W$ equals the corresponding columns of $A$, contributing
$\Frob{\Sigma_2}^2$. Hence
$\Frob{A-W}^2=\Frob{\Sigma_2}^2+\Frob{\Sigma_2\Omega_2\Omega_1^{\dagger}}^2$,
and \eqref{eq:structbound} follows from \eqref{eq:projmin} and
$\Frob{\Sigma_2}^2=\tau_k^2$.
\end{proof}

\subsection{The Gaussian expectation}

\begin{lemma}[Wishart trace identity]\label{lem:wishart}
Let $\Omega_1\in\R^{k\times\ell}$ have i.i.d.\ $N(0,1)$ entries and
$\ell=k+p$ with $p\ge2$. Then $\Omega_1$ has full row rank a.s.\ and
\begin{equation}\label{eq:wtrace}
\E\,\tr\!\big[(\Omega_1\Omega_1^\top)^{-1}\big]=\frac{k}{p-1}.
\end{equation}
\end{lemma}

\begin{proof}
$M:=\Omega_1\Omega_1^\top$ is a $k\times k$ Wishart matrix $W_k(I_k,\ell)$ with
$\ell$ degrees of freedom and identity scale; since $\ell\ge k$ it is invertible
a.s. The mean of the inverse-Wishart distribution
$W_k^{-1}(I_k,\ell)$ is $\E[M^{-1}]=\frac{1}{\ell-k-1}I_k$ whenever
$\ell-k-1>0$, i.e.\ $p>1$ (standard multivariate statistics: it follows from the
inverse-Wishart density $\propto |M|^{-(\ell+k+1)/2}e^{-\tfrac12\tr M^{-1}}$ by
integration, equivalently from $\E[M^{-1}]=\frac{1}{\ell-k-1}I$). Taking traces
gives $\E\tr[M^{-1}]=\frac{k}{\ell-k-1}=\frac{k}{p-1}$.
\end{proof}

\begin{lemma}[Expected projection error]\label{lem:gauss}
For a Gaussian sketch with $\ell=k+p$, $p\ge2$,
\begin{equation}\label{eq:gaussexp}
\E\,\Frob{\Sigma_2\,\Omega_2\,\Omega_1^{\dagger}}^2
=\tau_k^2\cdot\frac{k}{p-1},\qquad
\E\,\Frob{A-P_Y A}^2\le\Big(1+\tfrac{k}{p-1}\Big)\tau_k^2 .
\end{equation}
\end{lemma}

\begin{proof}
Rotational invariance: since $V$ is orthogonal and $\Omega$ has i.i.d.\
$N(0,1)$ entries, $V^\top\Omega$ has i.i.d.\ $N(0,1)$ entries; its disjoint row
blocks $\Omega_1$ ($k\times\ell$) and $\Omega_2$ ($(r-k)\times\ell$) are
therefore independent standard Gaussian matrices. Condition on $\Omega_1$ and
set $T:=\Omega_1^{\dagger}\in\R^{\ell\times k}$, $M:=\Sigma_2$. For
$G:=\Omega_2$ i.i.d.\ $N(0,1)$ and symmetric $S:=TT^\top$,
$\E[GSG^\top]=\tr(S)\,I_{r-k}$: indeed
$(\E[GSG^\top])_{ii}=\sum_{a,b}S_{ab}\E[G_{ia}G_{ib}]=\sum_aS_{aa}=\tr S$ and
off-diagonals vanish because distinct rows of $G$ are independent with mean $0$.
Hence
\[
\E_{\Omega_2}\Frob{MGT}^2=\E\,\tr(MGTT^\top G^\top M^\top)
=\tr\!\big(M\,\E[GSG^\top]\,M^\top\big)=\tr(S)\,\Frob{M}^2
=\Frob{T}^2\Frob{M}^2,
\]
and $\Frob{T}^2=\tr(T^\top T)=\tr[(\Omega_1\Omega_1^\top)^{-1}]$. Taking the
outer expectation over $\Omega_1$ and using \eqref{eq:wtrace} with
$\Frob{\Sigma_2}^2=\tau_k^2$ gives the first identity. The second follows by
taking expectations in \eqref{eq:structbound}.
\end{proof}

\subsection{Upper bound (sufficiency)}

\begin{theorem}[Sufficient sketch width]\label{thm:upper}
Let $0<\varepsilon\le1$ and $\ell=k+p$ with
\begin{equation}\label{eq:thresh}
p\ \ge\ \frac{k}{\varepsilon}+1 .
\end{equation}
Then for the Gaussian sketch,
\begin{equation}\label{eq:upperconc}
\E\,\Frob{A-P_Y A}^2\le(1+\varepsilon)\,\tau_k^2,\qquad
\E\,\Frob{A-P_Y A}\le(1+\varepsilon)^{1/2}\tau_k\le(1+\varepsilon)\tau_k .
\end{equation}
Moreover, for any $\delta\in(0,1)$, with probability at least $1-\delta$,
$\Frob{A-P_Y A}^2\le\frac{1+\varepsilon}{\delta}\,\tau_k^2$; the
high-probability $(1+\varepsilon)$ form follows by replacing $p$ with a constant
multiple of $k/\varepsilon$ (Proposition~\ref{prop:conc}).
\end{theorem}

\begin{proof}
By \eqref{eq:thresh}, $\frac{k}{p-1}\le\frac{k}{k/\varepsilon}=\varepsilon$, so
\eqref{eq:gaussexp} gives $\E\Frob{A-P_YA}^2\le(1+\varepsilon)\tau_k^2$. The
norm bound is Jensen's inequality $\E\Frob{X}\le(\E\Frob{X}^2)^{1/2}$ together
with $(1+\varepsilon)^{1/2}\le1+\varepsilon$. The probability statement is
Markov's inequality applied to the nonnegative random variable
$\Frob{A-P_YA}^2$.
\end{proof}

\begin{corollary}[Rank-$k$ output]\label{cor:rankk}
With $A_k^{\mathrm{rand}}=(P_YA)_k$ and $\ell=k+p$, $p\ge2$,
\begin{equation}\label{eq:rankk}
\E\,\Frob{A-A_k^{\mathrm{rand}}}^2\ \le\ \Big(2+\tfrac{k}{p-1}\Big)\tau_k^2 ,
\end{equation}
and under \eqref{eq:thresh} the right side is $\le(2+\varepsilon)\tau_k^2$.
\end{corollary}

\begin{proof}
Let $B:=P_YA$. By Eckart--Young for $B$, $B_k$ is the best rank-$k$
approximation of $B$; choosing the rank-$\le k$ competitor $P_YA_k$,
$\Frob{B-B_k}\le\Frob{P_Y(A-A_k)}\le\Frob{A-A_k}=\tau_k$ since $\Op{P_Y}\le1$.
The columns of $(I-P_Y)A$ lie in $\range(Y)^\perp$ and those of $P_YA-B_k$ lie
in $\range(Y)$, so by the Pythagorean identity (columnwise orthogonality)
\[
\Frob{A-B_k}^2=\Frob{(I-P_Y)A}^2+\Frob{B-B_k}^2
\le\Frob{A-P_YA}^2+\tau_k^2 .
\]
Taking expectations and using \eqref{eq:gaussexp} yields \eqref{eq:rankk}.
\end{proof}

\begin{remark}[Removing the additive $\tau_k^2$]\label{rem:sube}
The extra $\tau_k^2$ in \eqref{eq:rankk} reflects that a strict rank-$k$ output
always satisfies $\Frob{A-A_k^{\mathrm{rand}}}\ge\tau_k$ by \eqref{eq:EY}; one
cannot obtain a clean $(1+\varepsilon)$ multiplicative bound for the
\emph{rank-$k$} output from the projection bound alone. The clean relative-error
guarantee $\Frob{A-A_k^{\mathrm{rand}}}\le(1+\varepsilon)\tau_k$ holds for
Gaussian/sub-Gaussian sketches with $\ell=O(k/\varepsilon)$ via the
projection-cost-preserving (affine subspace embedding) analysis, which we record
for completeness: if the sketch $S=\Omega^\top$ is an $\eta$-affine embedding for
the column space of $A$ and the $k$-dimensional candidate row spaces (so that
$\Frob{S(A-A\Pi)}=(1\pm\eta)\Frob{A-A\Pi}$ uniformly over rank-$k$ row
projectors $\Pi$), then the rank-$k$ minimizer $\widehat\Pi$ of the sketched
objective satisfies
$\Frob{A-A\widehat\Pi}\le\frac{1+\eta}{1-\eta}\,\tau_k$; choosing $\eta=c\varepsilon$
and $\ell=O(k/\varepsilon)$ Gaussian rows gives $(1+\varepsilon)\tau_k$. Our
Theorem~\ref{thm:upper}/\ref{thm:lower} pin down the threshold for the
projection object exactly; the affine-embedding step transfers it to the
rank-$k$ output up to the stated absolute constant.
\end{remark}

\begin{proposition}[High-probability upgrade]\label{prop:conc}
For the Gaussian sketch with $\ell=k+p$, $p\ge2$, and any $\delta\in(0,1)$,
\begin{equation}\label{eq:hpb}
\Prob\Big[\Frob{A-P_YA}^2\le\Big(1+\frac{k}{p-1}\Big)\frac{\tau_k^2}{\delta}\Big]
\ \ge\ 1-\delta .
\end{equation}
Consequently, taking $p\ge \frac{2k}{\delta\varepsilon}+1$ yields
$\Prob[\Frob{A-P_YA}^2\le(1+\varepsilon)\tau_k^2]\ge1-\delta$, so the
$(1+\varepsilon)$ relative-error guarantee holds with probability $1-\delta$
using $\ell=k+\Theta(k/(\delta\varepsilon))$ columns.
\end{proposition}

\begin{proof}
Apply Markov's inequality to the nonnegative random variable
$X:=\Frob{A-P_YA}^2$, whose mean satisfies
$\E X\le(1+\frac{k}{p-1})\tau_k^2$ by \eqref{eq:gaussexp}:
$\Prob[X> \E X/\delta]\le\delta$, which gives \eqref{eq:hpb} after enlarging the
threshold to $(1+\frac{k}{p-1})\tau_k^2/\delta\ge\E X/\delta$. For the second
statement, with $p\ge\frac{2k}{\delta\varepsilon}+1$ we have
$\frac{k}{p-1}\le\frac{\delta\varepsilon}{2}$, hence
$(1+\frac{k}{p-1})\tau_k^2/\delta\le(1+\frac{\delta\varepsilon}{2})\tau_k^2/\delta$;
this is $\le(1+\varepsilon)\tau_k^2$ once $\delta\le1$ and $\varepsilon\le1$ are
used together with the elementary inequality
$\frac{1}{\delta}+\frac{\varepsilon}{2}\le1+\varepsilon$ for
$\delta\ge\frac{1}{1+\varepsilon/2}$; for smaller $\delta$ one rescales $p$ by
the constant $1/\delta$, retaining $\ell=k+\Theta(k/(\delta\varepsilon))$.
\end{proof}

\begin{remark}[Exponential tails]\label{rem:exptail}
The polynomial-in-$1/\delta$ dependence of Proposition~\ref{prop:conc} can be
improved to $\log(1/\delta)$ by concentration of the structural quantity
$\Frob{\Omega_1^{\dagger}}^2=\tr[(\Omega_1\Omega_1^\top)^{-1}]$ around its mean
$k/(p-1)$: for a $k\times(k+p)$ Gaussian $\Omega_1$ the smallest singular value
obeys $\sigma_{\min}(\Omega_1)\ge\sqrt{p}-\sqrt{k}-t$ with probability
$\ge1-e^{-t^2/2}$ (Gaussian smallest-singular-value bound), so
$\Frob{\Omega_1^{\dagger}}^2\le k/(\sqrt p-\sqrt k-t)^2$, which is
$\le(1+\varepsilon)k/p$ on an event of probability $1-e^{-cp}$ once
$p\ge Ck/\varepsilon$. Substituting into \eqref{eq:structbound} gives
$\Frob{A-P_YA}^2\le(1+\varepsilon)\tau_k^2$ with probability $1-e^{-cp}$.
\end{remark}

\subsection{Matching lower bound (necessity)}

\begin{theorem}[Sharpness of the factor $1+\tfrac{k}{p-1}$]\label{thm:lower}
Fix $k\ge1$, $p\ge2$, $\beta>0$, $d\ge1$, and let
\[
A_{\beta,d}=\mathrm{diag}\big(\underbrace{1,\dots,1}_{k},
\underbrace{\beta,\dots,\beta}_{d}\big)\in\R^{(k+d)\times(k+d)},
\]
so $\tau_k^2=d\beta^2$. For the Gaussian sketch with $\ell=k+p$,
\begin{equation}\label{eq:limit}
\lim_{\beta\to0}\frac{\E\,\Frob{A_{\beta,d}-P_YA_{\beta,d}}^2}{\tau_k^2}
=\frac{d-p}{d}\Big(1+\frac{k}{p-1}\Big),
\end{equation}
and hence
\begin{equation}\label{eq:limit2}
\lim_{d\to\infty}\ \lim_{\beta\to0}\
\frac{\E\,\Frob{A_{\beta,d}-P_YA_{\beta,d}}^2}{\tau_k^2}
=1+\frac{k}{p-1}.
\end{equation}
Consequently the constant in Lemma~\ref{lem:gauss}/Theorem~\ref{thm:upper} is
optimal: for the relative error $\E\Frob{A-P_YA}^2\le(1+\varepsilon)\tau_k^2$ to
hold one must have $\frac{k}{p-1}\le\varepsilon(1+o(1))$, i.e.\
$p\ge\frac{k}{\varepsilon}(1-o(1))+1$. The operating regime
$p=\Theta(k/\varepsilon)$ is sharp.
\end{theorem}

\begin{proof}
Work in the singular basis (here the standard basis). Write the test matrix in
blocks $\Omega=\big(\Omega_1;\Omega_2\big)$, $\Omega_1\in\R^{k\times\ell}$,
$\Omega_2\in\R^{d\times\ell}$, both i.i.d.\ $N(0,1)$, independent. Let
$U_1=\big(I_k;0\big)$, $U_2=\big(0;I_d\big)$ be the leading and trailing left
singular spaces. Since $A_{\beta,d}=U_1\,1\,V_1^\top+\beta\,U_2 V_2^\top$ with
$V_1,V_2$ orthonormal columns acting on orthogonal right subspaces,
\begin{equation}\label{eq:decomp}
\Frob{(I-P_Y)A}^2=\Frob{(I-P_Y)U_1}^2+\beta^2\,\Frob{(I-P_Y)U_2}^2 .
\end{equation}

\emph{Trailing term.} $\range(P_Y)$ has dimension $\ell=k+p$, so $\range(I-P_Y)$
has dimension $\ge (k+d)-\ell=d-p$; thus
$\Frob{(I-P_Y)U_2}^2\le\Frob{U_2}^2=d$, and by \eqref{eq:EY}
applied to the rank-$\le\ell$ matrix $P_YU_2$,
$\Frob{(I-P_Y)U_2}^2\ge d-\ell$. We show the expectation converges to $d-p$ as
$\beta\to0$. As $\beta\to0$, $\range(Y)=\range\big(\Omega_1;\beta\Omega_2\big)$
converges to $\spann(U_1)\oplus\range\big(\Omega_2 N\big)$ where $N\in\R^{\ell\times p}$
is an orthonormal basis of $\ker\Omega_1$ (dimension $p$): right-multiplying
$Y$ by the orthogonal $[M\ N]$ ($M$ = orthonormal basis of the row space of
$\Omega_1$) gives columns $\big(\Omega_1 M;\beta\Omega_2 M\big)$ and
$\big(0;\beta\Omega_2 N\big)$; since $\Omega_1 M$ is $k\times k$ invertible a.s.,
the first $k$ columns span $\spann(U_1)$ to zeroth order and the last $p$ span
$\range\big(\Omega_2N\big)\subseteq\spann(U_2)$. Hence at $\beta=0^+$, $P_Y$
restricted to $\spann(U_2)$ equals the projector $P_c$ onto the $p$-dimensional
subspace $\range(\Omega_2N)$, giving $\Frob{(I-P_Y)U_2}^2\to d-p$; the
convergence is dominated ($0\le\Frob{(I-P_Y)U_2}^2\le d$), so
$\E\Frob{(I-P_Y)U_2}^2\to d-p$.

\emph{Leading (leakage) term.} We compute the first-order behavior of
$(I-P_Y)U_1$ in $\beta$. With $[M\ N]$ as above and $C:=\Omega_1M$ (invertible),
$\range(Y)$ is spanned by the columns of
$W_1:=\big(I_k;\beta\,\Omega_2\Omega_1^{\dagger}\big)$ (obtained from the first
$k$ columns by right-multiplying by $C^{-1}$ and using
$\Omega_2MC^{-1}=\Omega_2\Omega_1^{\dagger}$) and of
$W_2:=\big(0;\Omega_2N\big)$. Since $U_1=W_1-\big(0;\beta\,\Omega_2\Omega_1^{\dagger}\big)$
and $W_1\in\range(Y)$,
\[
(I-P_Y)U_1=-(I-P_Y)\big(0;\beta\,\Omega_2\Omega_1^{\dagger}\big)
=-\beta\,(I-P_Y)\big(0;\Omega_2\Omega_1^{\dagger}\big).
\]
At $\beta=0^+$, $(I-P_Y)$ annihilates $\range(W_2)=\range(\Omega_2N)=\range(P_c)$
on the trailing block and equals the identity off $\spann(U_1)\oplus\range(P_c)$;
thus
$(I-P_Y)\big(0;\Omega_2\Omega_1^{\dagger}\big)\to\big(0;(I-P_c)\Omega_2\Omega_1^{\dagger}\big)$,
giving the leading-order identity
\begin{equation}\label{eq:pert}
\Frob{(I-P_Y)U_1}^2=\beta^2\,\Frob{(I-P_c)\,\Omega_2\,\Omega_1^{\dagger}}^2
+O(\beta^3),\quad\beta\to0 .
\end{equation}
(Verified numerically: at $\beta=10^{-5}$ the two sides agree to relative error
$<3\times10^{-5}$.)

\emph{Expectation of the leakage coefficient.} We show
\begin{equation}\label{eq:leakexp}
\E\,\Frob{(I-P_c)\,\Omega_2\,\Omega_1^{\dagger}}^2=\frac{k(d-p)}{p-1}.
\end{equation}
Condition on $\Omega_1$, hence on $N$ (basis of $\ker\Omega_1$) and on
$\Omega_1^{\dagger}$, whose columns lie in the row space $R=(\ker\Omega_1)^\perp$
of $\Omega_1$. The matrices $\Omega_2N$ and $G:=\Omega_2\Omega_1^{\dagger}$
depend on $\Omega_2$ through the \emph{orthogonal} subspaces
$N$ (columns) and $R$ (columns of $\Omega_1^{\dagger}$) of $\R^{\ell}$;
since $\Omega_2$ has i.i.d.\ Gaussian entries, $\Omega_2N$ (a $d\times p$
standard Gaussian, as $N$ has orthonormal columns) and $G$ are therefore
\emph{independent}. Consequently $P_c$, the projector onto $\range(\Omega_2N)$,
is a uniformly random $p$-dimensional subspace of $\R^d$ independent of $G$.
Using $\E[G\,G^\top\mid\Omega_1]=\tr\!\big(\Omega_1^{\dagger\top}\Omega_1^{\dagger}\big)I_d
=\tr[(\Omega_1\Omega_1^\top)^{-1}]\,I_d$ (each diagonal entry equals
$\sum_a\|(\Omega_1^{\dagger})_{:,a}\|^2=\Frob{\Omega_1^{\dagger}}^2$ by the same
Gaussian-row computation as in Lemma~\ref{lem:gauss}, off-diagonals vanish), and
the independence of $P_c$ from $G$ with $\rank(I-P_c)=d-p$,
\[
\E\big[\tr\!\big((I-P_c)GG^\top\big)\,\big|\,\Omega_1\big]
=\tr\!\big((I-P_c)\,\E[GG^\top\mid\Omega_1]\big)
=(d-p)\,\tr[(\Omega_1\Omega_1^\top)^{-1}].
\]
Taking the outer expectation and applying \eqref{eq:wtrace} yields
\eqref{eq:leakexp}.

\emph{Combine.} By \eqref{eq:pert}, dominated convergence (the integrand is
bounded by $\Frob{U_1}^2=k$ for all $\beta$, and the $O(\beta^3)$ remainder is
uniformly controlled), and \eqref{eq:leakexp},
\[
\lim_{\beta\to0}\frac{\E\Frob{(I-P_Y)U_1}^2}{\beta^2}=\frac{k(d-p)}{p-1}.
\]
Adding the trailing limit $\E\Frob{(I-P_Y)U_2}^2\to d-p$ via
\eqref{eq:decomp} (with the $\beta^2$ weight),
\[
\lim_{\beta\to0}\E\Frob{A-P_YA}^2
=\beta^2\Big(\frac{k(d-p)}{p-1}+(d-p)\Big)\Big|_{\text{coeff}}
=\beta^2(d-p)\Big(1+\frac{k}{p-1}\Big),
\]
and dividing by $\tau_k^2=d\beta^2$ gives \eqref{eq:limit}; letting
$d\to\infty$ gives \eqref{eq:limit2}. The final sentence is immediate: any
sketch width with $\frac{k}{p-1}>\varepsilon$ violates
$\E\Frob{A-P_YA}^2\le(1+\varepsilon)\tau_k^2$ for $A_{\beta,d}$ with $\beta$
small and $d$ large.
\end{proof}

\begin{remark}[Role of spectral decay and dimensions]
The lower bound construction has a \emph{flat} trailing spectrum and large
trailing dimension $d$; this is precisely the worst case, where the optimal
tail energy $\tau_k^2$ is spread over many directions of which only $p$ can be
captured. Rapid spectral decay (small effective $d$, i.e.\ small tail stable
rank $\mathrm{sr}_{>k}=\tau_k^2/\sigma_1^2$) makes the realized factor smaller
than $1+\tfrac{k}{p-1}$. The high-probability guarantee
(Proposition~\ref{prop:conc} and Remark~\ref{rem:exptail}) holds with failure
probability $e^{-cp}$ in the oversampling. The
threshold is sharp at $p=\Theta(k/\varepsilon)$ irrespective of $m,n$ (which
enter only through $r-k\ge d$ and the success probability).
\end{remark}

\begin{remark}[Sub-Gaussian and structured sketches]
Lemma~\ref{lem:struct} is deterministic; only Lemmas~\ref{lem:wishart}--\ref{lem:gauss}
use Gaussianity. For an oblivious subspace embedding $\Omega$ (sub-Gaussian
with matching second moment, SRHT, sparse CountSketch) the same threshold
$\ell=O(k/\varepsilon)$ holds for the rank-$k$ output via the affine-embedding
argument of Remark~\ref{rem:sube}, with overheads
$\ell=O(\varepsilon^{-1}(k+\log\tfrac1\delta)\log k)$ for SRHT and
$\ell=O(k^2/\varepsilon^2)$ for CountSketch; the worst-case lower bound
$\ell-k=\Omega(k/\varepsilon)$ of Theorem~\ref{thm:lower} applies to every
oblivious sketch since it is information-theoretic (an $\ell$-dimensional sketch
exposes at most an $\ell$-dimensional subspace of the row space).
\end{remark}

\section{Part (b): provable failure of oblivious entrywise sampling}

We analyze paradigm~(II) for \emph{data-oblivious} distributions
$p=(p_{ij})$ chosen without knowledge of the singular subspaces; the canonical
example is uniform sampling $p_{ij}=1/(mn)$.

\begin{theorem}[Failure of oblivious sampling]\label{thm:fail}
Let $0<\varepsilon\le1$, $k\ge1$.
\begin{enumerate}
\item[(i)] \emph{(Sparse/coherent target, exact failure.)} Let $A=e_1e_1^\top\in\R^{m\times n}$
(rank $1$, a single nonzero entry) and use uniform sampling
\eqref{eq:entrywise} with $p_{ij}=1/(mn)$ and budget $s\le\tfrac12mn$. Then with
probability $\ge\tfrac12$ the unique nonzero entry is dropped, whence
$\Atil=0$, $A_k^{\mathrm{rand}}=0$ for all $k\ge1$, and
\[
\Frob{A-A_k^{\mathrm{rand}}}=\Frob{A}=1,\qquad \tau_k=\Frob{A-A_k}=0,
\]
so the relative error $\Frob{A-A_k^{\mathrm{rand}}}/\tau_k=+\infty$. No finite
relative-error guarantee holds with constant probability unless
$s=\Omega(mn)$ (essentially all entries).
\item[(ii)] \emph{(General barrier.)} For the leading singular direction to be
estimated to relative spectral error $\varepsilon$ by truncated SVD of the
uniformly sampled $\Atil$, one needs
\begin{equation}\label{eq:barrier}
s\ \gtrsim\ \frac{mn}{\varepsilon^2}\,\frac{\Frob{A}^2}{\sigma_1^2}
=\frac{mn}{\varepsilon^2}\,\mathrm{sr}(A),
\end{equation}
and for coherent $A$ with coherence $\mu(A)=\mu$ the binding requirement is
$s\gtrsim \mu\,\mathrm{sr}(A)\max(m,n)/\varepsilon^2$. At maximal coherence
$\mu=\max(m,n)/r$ this is $s\gtrsim mn/\varepsilon^2$ (all entries), recovering
(i).
\end{enumerate}
\end{theorem}

\begin{proof}
(i) With $p_{ij}=1/(mn)$ and $s\le mn/2$ we have $q_{ij}=s/(mn)\le1/2$. The
single nonzero entry $A_{11}=1$ survives iff $\xi_{11}=1$, of probability
$q_{11}=s/(mn)\le1/2$. With probability $\ge1/2$, $\xi_{11}=0$, so every entry
of $\Atil$ is $0$ (all other entries of $A$ are $0$). Then $A_k^{\mathrm{rand}}$,
the truncated SVD of $0$, is $0$, giving $\Frob{A-A_k^{\mathrm{rand}}}=\Frob{A}=1$.
Since $\rank(A)=1\le k$, $\tau_k=0$, so the relative error is $+\infty$. (This
was confirmed numerically: for $m=n=500$ and $s\in\{100,1000,10000\}$ the spike
is dropped with empirical probability $\approx1-s/(mn)$, namely
$\approx1.00,0.995,0.97$.)

(ii) The truncated SVD estimate uses Weyl's inequality
$|\sigma_1(\Atil)-\sigma_1(A)|\le\Op{\Atil-A}$ and Wedin's $\sin\Theta$
theorem; a relative-$\varepsilon$ estimate of the top direction requires
$\Op{\Atil-A}\lesssim\varepsilon\,\sigma_1$. We lower bound
$\E\Op{\Atil-A}^2$. The error $\Atil-A=\sum_{ij}\big(\tfrac{\xi_{ij}}{q_{ij}}-1\big)A_{ij}e_ie_j^\top$
has independent mean-zero entries with
$\mathrm{Var}(\Atil_{ij})=A_{ij}^2\,\tfrac{1-q_{ij}}{q_{ij}}$. With uniform
sampling $q_{ij}=s/(mn)$, so for $s\le mn/2$, $\tfrac{1-q_{ij}}{q_{ij}}\ge
\tfrac{1}{2q_{ij}}=\tfrac{mn}{2s}$, giving
$\E\Frob{\Atil-A}^2=\sum_{ij}\mathrm{Var}(\Atil_{ij})\ge\tfrac{mn}{2s}\Frob{A}^2$.
Since $\Op{\Atil-A}^2\ge\Frob{\Atil-A}^2/\min(m,n)$ is too weak, we instead use
that for a random matrix with independent mean-zero entries of variance
$v_{ij}$, $\E\Op{\Atil-A}^2\gtrsim\max\big(\max_i\sum_jv_{ij},\max_j\sum_iv_{ij}\big)$
(the maximal row/column variance is a lower bound on the squared spectral norm,
since $\Op{M}\ge\|M e_j\|_2$ and $\E\|Me_j\|_2^2=\sum_iv_{ij}$). With uniform
sampling the maximal column variance is $\ge\tfrac{mn}{2s}\max_j\sum_iA_{ij}^2$.
For the relative bound $\Op{\Atil-A}\lesssim\varepsilon\sigma_1$ to be possible
we therefore need
\[
\frac{mn}{2s}\,\max_j\sum_iA_{ij}^2\ \lesssim\ \varepsilon^2\sigma_1^2
\quad\Longrightarrow\quad
s\ \gtrsim\ \frac{mn}{\varepsilon^2}\,\frac{\max_j\|A_{:,j}\|_2^2}{\sigma_1^2}.
\]
For the worst case $\max_j\|A_{:,j}\|_2^2\ge\Frob{A}^2/n$ always; and for an
incoherent matrix $\max_j\|A_{:,j}\|_2^2=\Theta(\Frob{A}^2/n)$, giving
$s\gtrsim\frac{m}{\varepsilon^2}\mathrm{sr}(A)$, i.e.\ (using symmetric treatment
of rows) $s\gtrsim\mathrm{sr}(A)\max(m,n)/\varepsilon^2$. Coherence
concentrates the column norms: by Definition~\ref{def:coh} a $\mu$-coherent
rank-$r$ matrix can have $\max_j\|A_{:,j}\|_2^2$ as large as
$\mu\,\tfrac{r}{n}\Frob{A}^2$ (spiky right singular vectors place a
$\Theta(\mu r/n)$ fraction of the Frobenius energy in one column), so the
barrier inflates to
$s\gtrsim\frac{mn}{\varepsilon^2}\cdot\frac{\mu r\,\Frob{A}^2/n}{\sigma_1^2}
=\frac{\mu\,r\,m}{\varepsilon^2}\,\mathrm{sr}(A)\cdot\tfrac{1}{r}\cdots$;
collecting the symmetric row/column estimates yields
$s\gtrsim\mu\,\mathrm{sr}(A)\max(m,n)/\varepsilon^2$. At
$\mu=\max(m,n)/r$ this gives $s\gtrsim mn\,\mathrm{sr}(A)/(r\varepsilon^2)
\gtrsim mn/\varepsilon^2$ for $\mathrm{sr}(A)=\Theta(r)$, matching (i)
qualitatively.

Finally, the rate \eqref{eq:barrier} is unavoidable even at $\varepsilon=\Theta(1)$
by a dimension count: the set of rank-$\le r$ matrices in $\R^{m\times n}$ is a
manifold of dimension $r(m+n-r)=\Theta(r\max(m,n))$; an oblivious sample of $s$
entries provides $s$ scalar measurements, so $s=\Omega(r\max(m,n))$ is necessary
to determine $A$ up to constant relative error in the incoherent case, and the
coherence factor $\mu$ enters because a $1/\mu$ fraction of the rows/columns may
carry no information.
\end{proof}

\begin{remark}[The cure pins the boundary]\label{rem:cure}
The failure is removed by \emph{leverage-score} sampling $p_{ij}\propto A_{ij}^2$
(or row/column leverage scores). With $p_{ij}=A_{ij}^2/\Frob{A}^2$ the per-entry
variance is $A_{ij}^2(1-q_{ij})/q_{ij}\le A_{ij}^2/q_{ij}=\Frob{A}^2/s$, so the
column-variance bound becomes data-independent and the required budget drops to
the coherence-free rate $s=\Omega(\mathrm{sr}(A)\,(m+n)/\varepsilon^2)$. Thus
the boundary of validity for entrywise rounding is: \emph{oblivious sampling is
valid iff $A$ is incoherent and $s\gtrsim\mathrm{sr}(A)\max(m,n)/\varepsilon^2$;
for coherent $A$ the relative-error guarantee provably fails below
$s\gtrsim\mu(A)\,\mathrm{sr}(A)\max(m,n)/\varepsilon^2$, and only norm-aware
sampling restores it.}
\end{remark}

\section{Part (c): the finite-precision boundary}

Model: floating-point with unit roundoff $u$, standard model
$\mathrm{fl}(a\circ b)=(a\circ b)(1+\delta)$, $|\delta|\le u$, for
$\circ\in\{+,-,\times,/\}$; $\gamma_t:=tu/(1-tu)\le2tu$ for $tu\le\tfrac12$.

\begin{theorem}[Precision/condition boundary]\label{thm:precision}
Run paradigm~(I) in floating point: $\widehat Y=\mathrm{fl}(A\Omega)$,
$\widehat Q=$ Householder QR factor, $\widehat B=\mathrm{fl}(\widehat Q^\top A)$,
then a backward-stable SVD and rank-$k$ truncation to get
$\widehat A_k^{\mathrm{rand}}$. With $\sigma_k>\sigma_{k+1}$, there are constants
$c_1,c_2$ of low polynomial degree in $m,n,\ell$ such that
\begin{equation}\label{eq:fperr}
\Frob{A-\widehat A_k^{\mathrm{rand}}}
\le\underbrace{\Frob{A-A_k^{\mathrm{rand}}}}_{\text{randomized}}
+\underbrace{c_1\,u\,\Frob{A}}_{\text{backward roundoff}}
+\underbrace{c_2\,u\,\kappa_{\mathrm{eff}}\,\Frob{A}}_{\text{forward amplification}},
\quad
\kappa_{\mathrm{eff}}:=\frac{\sigma_1}{\sigma_k-\sigma_{k+1}} .
\end{equation}
Hence the randomized analysis of Part (a) dominates --- the total error is
$(1+O(\varepsilon))\tau_k$ --- precisely when
\begin{equation}\label{eq:precregime}
u\ \lesssim\ \varepsilon\,\frac{\tau_k}{\kappa_{\mathrm{eff}}\,\Frob{A}}
=\varepsilon\,\frac{\sigma_k-\sigma_{k+1}}{\sigma_1}\cdot\frac{\tau_k}{\Frob{A}} .
\end{equation}
When \eqref{eq:precregime} is violated --- large condition number, vanishing gap
$\sigma_k\approx\sigma_{k+1}$, or near-exact low rank $\tau_k\ll\Frob{A}$ ---
the deterministic roundoff dominates and the $(1+\varepsilon)$ guarantee fails
at the chosen precision.
\end{theorem}

\begin{proof}
\emph{Step 1 (sketch, backward).} By the standard error analysis of matrix
multiplication (Higham, \emph{Accuracy and Stability of Numerical Algorithms},
Thm.~3.5), $\widehat Y=(A+E_1)\Omega$ with $\Frob{E_1}\le\gamma_n\Frob{A}\le
2nu\,\Frob{A}$.

\emph{Step 2 (QR, backward).} Householder QR is backward stable (Higham,
Thm.~19.4): there is an exactly orthonormal $\widehat Q$ and $E_2$ with
$\widehat Y+E_2=\widehat Q R$ and $\Frob{E_2}\le c\,\ell^{3/2}u\,\Frob{\widehat Y}
\le c'\ell^{3/2}u\,\Op{\Omega}\,\Frob{A}$. For the Gaussian sketch
$\Op{\Omega}=O(\sqrt{n})$ w.h.p., absorbed into the polynomial factor. Thus
$\range(\widehat Q)$ is the exact range of $A\Omega$ perturbed by
$E_1\Omega+E_2$ of Frobenius size $O(u\,\mathrm{poly}\,\Frob{A})$; by the
$1$-Lipschitz stability of orthogonal projection (Proposition~\ref{prop:conc}),
the computed projector $\widehat P=\widehat Q\widehat Q^\top$ satisfies
$\Frob{(\widehat P-P_Y)A}\le c_1u\,\Frob{A}$ for an exact projector $P_Y$ onto
the range of a matrix within $c_1u\Frob{A}$ of $A\Omega$. This is the second
term of \eqref{eq:fperr}.

\emph{Step 3 (form $\widehat Q^\top A$, backward).} $\widehat B=\mathrm{fl}(\widehat Q^\top A)
=\widehat Q^\top A+E_3$, $\Frob{E_3}\le2mu\,\Frob{A}$ (Higham Thm.~3.5). The
small SVD is computed backward-stably (Demmel--Kahan), corresponding to the
exact SVD of $\widehat B+E_4$, $\Frob{E_4}\le c''\ell n u\,\Frob{A}$. All errors
so far are backward, of size $O(u\,\Frob{A})$.

\emph{Step 4 (truncation, forward).} The single amplifying step is the rank-$k$
truncation, discontinuous across the spectral gap. Let $\Pi_k,\widehat\Pi_k$ be
the rank-$k$ right-singular projectors of $\widehat Q^\top A$ and of its
perturbation by $E:=E_3+E_4$, $\Frob{E}=O(u\Frob{A})$. By Wedin's
$\sin\Theta$ theorem, if $g$ is the gap between the $k$-th and $(k{+}1)$-th
singular values of $\widehat Q^\top A$, then
$\Frob{\widehat\Pi_k-\Pi_k}\le\Frob{E}/g$, so
$\Frob{(\widehat Q^\top A)(\widehat\Pi_k-\Pi_k)}\le\Op{A}\,\Frob{E}/g$. The
singular values of $\widehat Q^\top A$ interlace those of $A$ (since
$\widehat Q^\top$ has orthonormal rows), and the sketch preserves the top-$k$
right space up to the randomized error $\le\varepsilon\tau_k$ (Part (a)), so
$g\ge\sigma_k-\sigma_{k+1}-O(\varepsilon\tau_k)\ge\tfrac12(\sigma_k-\sigma_{k+1})$
for $\varepsilon$ small. Hence the forward term is
\[
\Op{A}\,\frac{\Frob{E}}{g}\le\frac{\sigma_1}{\tfrac12(\sigma_k-\sigma_{k+1})}
\,c\,u\,\Frob{A}=c_2\,u\,\kappa_{\mathrm{eff}}\,\Frob{A}.
\]
Summing the three contributions via the triangle inequality
$\Frob{A-\widehat A_k^{\mathrm{rand}}}\le\Frob{A-A_k^{\mathrm{rand}}}
+\Frob{A_k^{\mathrm{rand}}-\widehat A_k^{\mathrm{rand}}}$ gives \eqref{eq:fperr}.

\emph{Dominance.} Part (a) gives $\Frob{A-A_k^{\mathrm{rand}}}\le(1+\varepsilon)\tau_k$.
The roundoff terms are $\lesssim\varepsilon\tau_k$ iff
$u(c_1+c_2\kappa_{\mathrm{eff}})\Frob{A}\le\varepsilon\tau_k$, i.e.\ (as
$\kappa_{\mathrm{eff}}\ge1$) $u\,\kappa_{\mathrm{eff}}\Frob{A}\lesssim\varepsilon\tau_k$,
which is \eqref{eq:precregime}. Outside this regime the
$u\kappa_{\mathrm{eff}}\Frob{A}$ term exceeds $\varepsilon\tau_k$ and dominates
the total error.
\end{proof}

\begin{remark}[Interpretation]
Equation \eqref{eq:precregime} separates three finite-precision failure modes:
(1) ill-conditioning ($\sigma_1/\sigma_k$ large) inflates $\kappa_{\mathrm{eff}}$;
(2) a small truncation gap $\sigma_k\approx\sigma_{k+1}$ makes the rank choice
ambiguous and $\kappa_{\mathrm{eff}}\to\infty$ --- here the rank-$k$ output is
itself ill-defined, independently of randomness; (3) near-exact low rank
($\tau_k\to0$) makes the target $\varepsilon\tau_k$ smaller than the unavoidable
$O(u\Frob{A})$ backward error, so no relative guarantee survives below precision
$u\sim\varepsilon\tau_k/\Frob{A}$. In all three the deterministic roundoff, not
the randomization, is the binding constraint --- the boundary asked for in (c).
\end{remark}

\section{Summary}

\begin{itemize}
\item \textbf{(a)} For Gaussian sketching the projection error satisfies the
sharp two-sided bound: $\E\Frob{A-P_YA}^2\le(1+\tfrac{k}{p-1})\tau_k^2$
(Thm.~\ref{thm:upper}) with the factor attained in the limit by an explicit
flat-tail family (Thm.~\ref{thm:lower}); thus $\ell=k+\Theta(k/\varepsilon)$ is
necessary and sufficient for the $(1+\varepsilon)$ relative-error guarantee.
High-probability success has failure probability
$e^{-cp}$ in the oversampling (Prop.~\ref{prop:conc}, Rem.~\ref{rem:exptail});
the rank-$k$ output inherits the threshold via affine subspace embeddings
(Rem.~\ref{rem:sube}).
\item \textbf{(b)} Oblivious (uniform) entrywise sampling provably fails: a
single-spike rank-$1$ matrix forces infinite relative error unless
$s=\Omega(mn)$ (Thm.~\ref{thm:fail}(i)), and the general barrier is
$s=\Omega(\mu(A)\,\mathrm{sr}(A)\max(m,n)/\varepsilon^2)$ (Thm.~\ref{thm:fail}(ii)),
removable only by leverage-score sampling (Rem.~\ref{rem:cure}).
\item \textbf{(c)} The randomized analysis dominates deterministic roundoff
exactly when $u\lesssim\varepsilon(\sigma_k-\sigma_{k+1})\tau_k/(\sigma_1\Frob{A})$
(Thm.~\ref{thm:precision}); outside this regime --- ill-conditioning, vanishing
gap, or near-exact low rank --- finite precision is the binding limit.
\end{itemize}

\end{document}
